本文将复Hamiltonian表示、保守生成函数中点构造及式（36）的数据生成联系到实复学习与残差补偿。非单位标量、构型依赖质量及耦合矩阵基准均观察到近二阶收敛；保守分支的首步辛缺陷接近浮点精度，阻尼中点的功率平衡则不等于连续时变Birkhoff二形式保持。

四组模型的五项配对、物理基线及步长加密表明，所测弱非线性区间内已知物理主部解释主要精度收益。构型依赖质量坐标中重新生成的非零$\Delta$数据也支持这一结论；复值网络较低的教师拟合误差没有形成稳定的连续预测优势，杆端与完整状态须分别判断。该坐标检验不等于实体质量本构标定。

同梁磁势提供领先阶非退化性，连续频率与步长频差揭示其分辨尺度。域内$p=6$梯度量和能量包络进一步量化学习、物理扰动及离散误差预算；解析或高阶范数、全阶非共振与阈值常数仍待认证，现有结果不构成严格KAM环面证书。

残差补偿在独立测试状态上降低了力误差；初值匹配诊断中，两类学习补偿在两档步长均改善杆端与状态指标，而共同零参考下的收益较弱。有势外生输入下I/G与对称组合约二阶，步后加载趋一阶；实际采样反馈相对连续反馈的总误差仍近一阶，完整反馈映射也未保辛。本文据此支持在明确质量维度、控制可实现方向和参考定义后，分别评价精度、补偿收益与结构。一般耗散复生成量、完整四场及高模态效应和实物闭环仍需独立验证。
